\documentclass[coursework]{SCWorks}
% Тип обучения (одно из значений):
%    bachelor   - бакалавриат (по умолчанию)
%    spec       - специальность
%    master     - магистратура
% Форма обучения (одно из значений):
%    och        - очное (по умолчанию)
%    zaoch      - заочное
% Тип работы (одно из значений):
%    coursework - курсовая работа (по умолчанию)
%    referat    - реферат
%  * otchet     - универсальный отчет
%  * nirjournal - журнал НИР
%  * digital    - итоговая работа для цифровой кафедры
%    diploma    - дипломная работа
%    pract      - отчет о научно-исследовательской работе
%    autoref    - автореферат выпускной работы
%    assignment - задание на выпускную квалификационную работу
%    review     - отзыв руководителя
%    critique   - рецензия на выпускную работу

\usepackage{preamble}
\usepackage{xurl}

\begin{document}

% Кафедра (в родительном падеже)
\chair{математической кибернетики и компьютерных наук}

% Тема работы
\title{Архитектура образовательной платформы: интеграция Git для версионного контроля и приема решений}

% Курс
\course{2}

% Группа
\group{251}

% Факультет (в родительном падеже) (по умолчанию "факультета КНиИТ")
% \department{факультета КНиИТ}

% Специальность/направление код - наименование
% \napravlenie{02.03.02 "--- Фундаментальная информатика и информационные технологии}
% \napravlenie{02.03.01 "--- Математическое обеспечение и администрирование информационных систем}
% \napravlenie{09.03.01 "--- Информатика и вычислительная техника}
\napravlenie{09.03.04 "--- Программная инженерия}
% \napravlenie{10.05.01 "--- Компьютерная безопасность}

% Для студентки. Для работы студента следующая команда не нужна.
% \studenttitle{студентки}

% Фамилия, имя, отчество в родительном падеже
\author{Железко Александра Дмитриевича}

% Заведующий кафедрой 
\chtitle{доцент, к.\,ф.-м.\,н.}
\chname{С.\,В.\,Миронов}

% Руководитель ДПП ПП для цифровой кафедры (перекрывает заведующего кафедры)
% \chpretitle{
%     заведующий кафедрой математических основ информатики и олимпиадного\\
%     программирования на базе МАОУ <<Ф"=Т лицей №1>>
% }
% \chtitle{г. Саратов, к.\,ф.-м.\,н., доцент}
% \chname{Кондратова\, Ю.\,Н.}

% Научный руководитель (для реферата преподаватель проверяющий работу)
\satitle{ассистент} %должность, степень, звание
\saname{Н.\,С.\,Рыданов}

% Руководитель практики от организации (руководитель для цифровой кафедры)
\patitle{доцент, к.\,ф.-м.\,н.}
\paname{С.\,В.\,Миронов}

% Руководитель НИР
\nirtitle{доцент, к.\,п.\,н.} % степень, звание
\nirname{В.\,А.\,Векслер}

% Семестр (только для практики, для остальных типов работ не используется)
\term{2}

% Наименование практики (только для практики, для остальных типов работ не
% используется)
\practtype{учебная}

% Продолжительность практики (количество недель) (только для практики, для
% остальных типов работ не используется)
\duration{2}

% Даты начала и окончания практики (только для практики, для остальных типов
% работ не используется)
\practStart{01.07.2024}
\practFinish{13.01.2024}

% Год выполнения отчета
\date{2026}

\maketitle

% Включение нумерации рисунков, формул и таблиц по разделам (по умолчанию -
% нумерация сквозная) (допускается оба вида нумерации)
\secNumbering

\tableofcontents

% Раздел "Обозначения и сокращения". Может отсутствовать в работе
% \abbreviations
% \begin{description}
%     \item ... "--- ...
%     \item ... "--- ...
% \end{description}

% Раздел "Определения". Может отсутствовать в работе
% \definitions

% Раздел "Определения, обозначения и сокращения". Может отсутствовать в работе.
% Если присутствует, то заменяет собой разделы "Обозначения и сокращения" и
% "Определения"
% \defabbr

\intro

Современное высшее образование, особенно в области информационных технологий, находится в условиях активной трансформации. Вызвано это как нехваткой квалифицированных кадров на рынке труда, так и ростом числа желающих развиваться в этой сфере. На фоне развития (и локального регресса) отрасли, традиционные образовательные подходы и методики обучения теряют свою актуальность, становятся недостаточно эффективными в подготовке специалистов. Прежние стандартные системы не всегда способны обеспечить необходимую гибкость, интерактивность и практическую направленность, которые востребованы в условиях быстро меняющегося рынка.

Под влиянием данных факторов активно развиваются цифровые образовательные платформы. Они не только помогают в организации и структурировании учебных программ, но и облегчают повседневную работу как студентов, так и преподавателей, автоматизируя рутинные процессы и предоставляя современные инструменты для обучения. Одной из ключевых задач таких платформ --- организация процесса сдачи, хранения и проверки заданий.

Существующие подходы к управлению студенческими работами, которые основаные на сохранении файлов в виде BLOB-объектов в базах данных или на файловых системах, имеют ряд недостатков: неэффективное хранение, отсутствие механизмов версионирования, затруднение с обеспечением целостности данных. Эти ограничения приводят к трудностям при реализации функционала сравнения версий (diff), восстановлении предыдущих состояний работ студентов и управлении метаданными каждой попытки сдачи, что снижает общую эффективность образовательного процесса.

В индустрии IT для решения проблем версионирования и хранения проектов, а также совместной работы над ними используются системы контроля версий как, например, Git. Тут нужно что-то продолжить про полезность навыка использования git, и что использование этого инструмента повседневно в целях обучения может сильно помочь в будущем. Просто нужно правильно сформулировать.

Целью данной работы является проектирование и реализация части образовательной платформы, которая использует Git для версионирования студенческих заданий, эффективного хранения данных, а также функционалом для удобной сдачи и проверки решений.

В ходе работы должны быть решены следующие задачи:
\begin{itemize}
	\item проанализировать существующие подходы к хранению программного кода, выявив преимущества систем контроля версий на базе Git;
	\item исследовать применение архитектуры REST API для построения основного сервера образовательной платформы;
	\item изучить принципы работы протокола SSH в обеспечении безопасного транспортного уровня для Git;
	\item разработать архитектуру Git-сервера;
	\item реализовать механизм аутентификации и авторизации с использованием SSH-ключей, а также управление доступом к репозиториям;
	\item внедрить механизм передачи метаданных в процесс git push для управления логикой приёма заданий;
	\item создать комбинированную модель хранения попыток, совмещающую Git для файлов и реляционную базу данных для метаданных.
\end{itemize}

\conclusion

% Окончание основного документа и начало приложений Каждая последующая секция
% документа будет являться приложением
\appendix

\end{document}

